%%%%%%%%%%%%%%%%%%%%%%%%%%%%%%%%%%%%%%%%%%%%%%%%%%%%%%%%%%%%%%%%%%%%
%% I, the copyright holder of this work, release this work into the
%% public domain. This applies worldwide. In some countries this may
%% not be legally possible; if so: I grant anyone the right to use
%% this work for any purpose, without any conditions, unless such
%% conditions are required by law.
%%%%%%%%%%%%%%%%%%%%%%%%%%%%%%%%%%%%%%%%%%%%%%%%%%%%%%%%%%%%%%%%%%%%

\documentclass{beamer}
\usetheme[faculty=ped,basePath=../fibeamer,logo=../fibeamer/logo/mu/Logo-UC-3.png]{fibeamer}
\graphicspath{{../resources/}}

\usepackage{algorithm}
\usepackage{algpseudocode}

\usetheme[faculty=ped]{fibeamer}
\usepackage[utf8]{inputenc}
\usepackage[
  main=english, %% By using `czech` or `slovak` as the main locale
                %% instead of `english`, you can typeset the
                %% presentation in either Czech or Slovak,
                %% respectively.
  czech, slovak %% The additional keys allow foreign texts to be
]{babel}        %% typeset as follows:
%%
%%   \begin{otherlanguage}{czech}   ... \end{otherlanguage}
%%   \begin{otherlanguage}{slovak}  ... \end{otherlanguage}
%%
%% These macros specify information about the presentation
\title{Control de flujo: Soluciones propuestas} %% that will be typeset on the
\subtitle{Semestre II-2026} %% title page.
\author{Andreina Cota Vidaurre}
%% These additional packages are used within the document:
\usepackage{ragged2e}  % `\justifying` text
\usepackage{booktabs}  % Tables
\usepackage{tabularx}
\usepackage{tikz}      % Diagrams
\usetikzlibrary{calc, shapes, backgrounds}
\usetikzlibrary{positioning, arrows.meta, shadows.blur, shapes.geometric}
\usepackage{amsmath, amssymb}
\usepackage{url}       % `\url`s
\usepackage{listings}  % Code listings
\usepackage{graphicx}
\usepackage{amsmath}
\usepackage[T1]{fontenc}
\usepackage{xcolor}
\usepackage{minted}
% \usepackage{adjustbox}


% Definicion de pseudocodigo en espaniol
\lstdefinelanguage{PseudocodigoES}{
    morekeywords={INICIO,FIN,PARA,HACER,SI,ENTONCES,FIN_SI,FIN_PARA,IMPRIMIR},
    sensitive=true,
    morecomment=[l]{//},
    morestring=[b]"
}

% Definicion de pseudocodigo en ingles
\lstdefinelanguage{PseudocodeEN}{
    morekeywords={BEGIN,END,FOR,DO,IF,THEN,END_IF,END_FOR,PRINT},
    sensitive=true,
    morecomment=[l]{//},
    morestring=[b]"
}

\lstset{
    basicstyle=\ttfamily\small,
    keywordstyle=\bfseries,
    columns=fullflexible,
    keepspaces=true,
    showstringspaces=false,
    frame=single,
    xleftmargin=0em,
    framexleftmargin=0em,
    framesep=2pt,
    gobble=4,
    resetmargins=true,
    aboveskip=0pt,
    belowskip=0pt
}

\lstdefinestyle{pythonstyle}{
    language=Python,
    basicstyle=\ttfamily\small,
    keywordstyle=\bfseries,
    commentstyle=\itshape,
    stringstyle=\ttfamily,
    showstringspaces=false,
    columns=fullflexible,
    keepspaces=true,
    frame=single,
    breaklines=true,
    xleftmargin=0em,
    framexleftmargin=0em,
    framesep=2pt,
    aboveskip=0pt,
    belowskip=0pt,
    gobble=4,
    resetmargins=true,
    emph={print,input,int,float,str,bool,len,range},
    emphstyle=\bfseries
}

\lstset{style=pythonstyle}

% Estilos del diagrama
\tikzstyle{iniciofin} = [
    ellipse,
    minimum width=2.5cm,
    minimum height=1cm,
    text centered,
    draw=black,
    fill=blue!25
]

\tikzstyle{proceso} = [
    rectangle,
    rounded corners,
    minimum width=3cm,
    minimum height=1cm,
    text centered,
    draw=black,
    fill=cyan!20
]

\tikzstyle{decision} = [
    diamond,
    aspect=2,
    text centered,
    draw=black,
    fill=yellow!35
]

\tikzstyle{flecha} = [
    thick,
    ->,
    >=Stealth
]

% \definecolor{green_python}{HTML}{52A736}

% \newcolumntype{Y}{>{\raggedright\arraybackslash}X}

\frenchspacing
\begin{document}
  \shorthandoff{-}
  \frame[c]{\maketitle}

% \begin{darkframes}

    \begin{frame}[fragile,t]{Ejercicio práctico}
        En un servicio de urgencias, un paciente se clasifica como
        \textbf{prioridad alta} si su presión sistólica es menor a
        90\,mmHg \textbf{o} su frecuencia respiratoria supera las
        30 respiraciones por minuto.
        \begin{minted}[xleftmargin=-2cm, frame=leftline, fontsize=\small]{python}
            presion_sistolica = 85      # mmHg
            frecuencia_respiratoria = 22  # resp/min
            
            # Completa la condicion:
            if presion_sistolica < 90 or frecuencia_respiratoria > 30:
                print("Prioridad alta")
            else:
                print("Prioridad estandar")
        \end{minted}
    \end{frame}

    \begin{frame}{Alerta por calidad de aire - Ejercicio}
        Dado el nivel de PM2.5 en el aire ($\mu$g/m³) y si la persona consultante pertenece 
        a un grupo de riesgo (niños, adultos mayores o personas con enfermedades respiratorias), 
        determina el mensaje de alerta que debe mostrar la aplicación.

        \begin{itemize}
            \item Si el PM2.5 supera 55 $\mu$g/m³ (nivel considerado dañino según estándares de 
            calidad del aire), y la persona es de grupo de riesgo → alerta máxima.
            \item Si supera ese nivel pero la persona no es de grupo de riesgo $\rightarrow$ alerta moderada.
            \item Si no supera ese nivel $\rightarrow$ aire aceptable.
        \end{itemize}
    \end{frame}

    \begin{frame}[fragile,t]{Alerta por calidad de aire - Solución propuesta}
        \begin{minted}[xleftmargin=-2cm, frame=leftline, fontsize=\small]{python}
            pm25 = 62              # microgramos por metro cubico
            grupo_riesgo = True

            if pm25 > 55:
                if grupo_riesgo:
                    print("Alerta maxima: evitar salir")
                else:
                    print("Alerta: reducir actividades al aire libre")
            else:
                print("Aire aceptable")
        \end{minted}
    \end{frame}

    \begin{frame}{Autorización de vuelo - Ejercicio}
        Una empresa de delivery con drones te pide programar la decisión de si un dron 
        puede iniciar su ruta. Los drones no pueden volar con vientos superiores a 40 km/h 
        por razones de estabilidad. Si el viento lo permite, el vuelo depende de la batería: 
        se requiere al menos 30\% para cualquier entrega, pero si el pedido pesa más de 2 kg, 
        el consumo aumenta y se exige un mínimo de 50\%.

        A partir del viento, el nivel de batería y el peso del pedido, determina si el dron 
        puede iniciar la entrega y, si no puede, indica el motivo.
    \end{frame}

    \begin{frame}[fragile,t]{Autorización de vuelo - Solución propuesta}
        \begin{minted}[xleftmargin=-2cm, frame=leftline, fontsize=\tiny]{python}
                viento = 25          # km/h
                bateria = 45          # %
                peso_pedido = 2.5     # kg

                if viento > 40:
                    print("No autorizado: viento excesivo")
                else:
                    if peso_pedido > 2:
                        if bateria >= 50:
                            print("Autorizado")
                        else:
                            print("No autorizado: bateria insuficiente para el peso")
                    else:
                        if bateria >= 30:
                            print("Autorizado")
                        else:
                            print("No autorizado: batería insuficiente")
        \end{minted}
    \end{frame}

    \begin{frame}{Clasificación de hipotermia - Ejercicio}
        Trabajas en una app para equipos de rescate en montaña. El protocolo clínico distingue tres 
        niveles de hipotermia según la temperatura corporal central: leve (entre 32\,°C y 35\,°C), 
        moderada (entre 28\,°C y 32\,°C) y severa (bajo 28\,°C). El equipo médico te indica que, además, 
        si el paciente ha estado expuesto al frío por más de 3 horas, su caso debe marcarse como de 
        mayor urgencia dentro de la categoría que le corresponda.

        A partir de la temperatura corporal y el tiempo de exposición de un paciente, escribe un 
        programa que entregue un mensaje con su nivel de hipotermia y su nivel de urgencia.
    \end{frame}

    \begin{frame}[fragile,t]{Clasificación de hipotermia - Solución propuesta}
        \begin{minted}[xleftmargin=-2cm, frame=leftline, fontsize=\tiny]{python}
                temperatura_corporal = 30.5   # grados Celsius
                horas_exposicion = 4

                if temperatura_corporal < 28:
                    if horas_exposicion > 3:
                        print("Hipotermia severa - urgencia critica")
                    else:
                        print("Hipotermia severa - urgencia alta")
                elif temperatura_corporal < 32:
                    if horas_exposicion > 3:
                        print("Hipotermia moderada - urgencia alta")
                    else:
                        print("Hipotermia moderada - urgencia media")
                else:
                    if horas_exposicion > 3:
                        print("Hipotermia leve - urgencia media")
                    else:
                        print("Hipotermia leve - urgencia baja")
        \end{minted}
    \end{frame}
    
\end{document}
